\section{The lesion}
\label{sec:lesion}
Before deriving a mechanism we locate the damage among the 224 quantized matrices, excise it and transplant it (\cref{tab:lesion,tab:transplant}); we call its concentrated part the \emph{lesion}.

\paragraph{Where it is.}
Comparing hidden states of the collapsed and the full-precision Llama singles out one position and one depth: at the output of layer 2 the relative error at position \LlesionPos{}, the token \texttt{<|end\_header\_id|>} that closes the chat header, is \LposErr{} against \LposErrRest{} elsewhere; for Qwen it is \QsinkTok{} at the output of layer 5 (\QposErr{}).
These are the depths at which the attention sink forms and the positions that act as secondary sinks, and the matrix that writes the sink state is the MLP down-projection of the sink-forming layer, layer 1 for Llama and layer \QlesionLayer{} for Qwen.
Re-rounding that matrix without compensation removes what the compensation wrote: with the other 223 matrices untouched, Llama rises from \LgptqIF{} to \LoneMatrixIF{} and Qwen from \QgptqIF{} to \QoneMatrixIF{}, while the attention projections of the first three layers (\LearlyAttnIF{}) or the five largest rows of the matrix (\LswRowsIF{}) at full precision do not help: not an attention defect, not the super weights' magnitude, but what the compensation wrote around them.

\paragraph{Necessary, not sufficient.}
Copied into the full-precision Llama the lesion costs almost nothing (\TPfp{}), into the RTN network four points (\TPrtn{}), and Qwen's costs the full-precision Qwen \TPqfpDrop{} points (\TPqfp{}), short of collapse: one matrix does not collapse a model on its own.
Inside a GPTQ-quantized network it does, and not only its own: copied into the healthy window-calibrated network (\LcwinIF{}) it collapses it (\TPcwin{}), and the window checkpoint's matrix, which its own network tolerates, does not cure the collapsed short-document network (\TPshort{}) whereas the RTN version of the same matrix does (\LoneMatrixIF{}).
\paragraph{What else it needs.}
Not the layers quantized after it: rounding every matrix from layer $k$ onward to nearest, for $k$ from 2 to 24, leaves the collapse intact (\RtnTailRange{}), whereas the same lesion inside an all-RTN network scores \TPrtn{}.
Transplanting subsets of the first two layers into the RTN network (\cref{tab:sets}; single draws, seed noise \SeedNoise{} points) isolates the one thing it needs, the down-projection of layer 0: the two down-projections together collapse the RTN network (\TwoDownIF{}), as does every tested set containing them, including all fourteen matrices of the two layers (\LfourteenIF{}), and every tested set without the layer-0 down-projection is healthy, among them the whole of layer 1 with the lesion inside it (\LoneWholeIF{}).
This pair is Llama's minimal collapsing set; each matrix is necessary, since removing either cures (\LrtnZeroDownIF{}, \LoneMatrixIF{}) and neither alone does anything (\LzeroAloneIF{}, \LoneAloneIF{} against \LrtnIF{}).
The layer-0 down-projection is a weaker copy of the same object (\bos{} share of its Hessian trace \AbosZero{} against \AbosOne{} at layer 1 and at most \AbosRest{} elsewhere; overshoot and row concentration in \cref{app:extra}), and \bos{}-heavy only under short documents, where a \bos{} token arrives every few hundred tokens: its share is \AbosZeroWin{} under C4 windows and \AbosZeroPile{} under Pile windows, the protocols Llama survives.
Only excision tests this, since a foreign GPTQ matrix is fatal in any GPTQ background: the window version of the layer-0 matrix leaves the collapse (\LcwinZeroSwapIF{}), the window versions of both cure it (\LcwinPairSwapIF{}).
The set is model-dependent: on Qwen the sink-forming down-projection alone collapses the RTN network (\QlesionRtnIF{}, below the fully quantized network's \QgptqIF{}), the second matrix adds nothing and excising both cures (\QtwoMatrixIF{}); on Mistral the window collapse needs the fourteen matrices of layers 0 and 1 (\MfourteenIF{}) but neither the down-projections (\MfourDownIF{}) nor layer 1 alone (\MlayerOneIF{}), a phenotype we localize but do not resolve.
\paragraph{Fatal wherever nothing absorbs it; the neighbors decide.}
The calibration draw that leaves Llama healthy (\LgptqSeedIF{}) writes a pair as fatal to the RTN network as the collapsing draw's (\LseedOnePairIF{} against \TwoDownIF{}), but its fourteen matrices of layers 0 and 1 together score \LseedOneFourteenIF{}, its own value, where the collapsing draw's score \LseedZeroFourteenIF{}; Tulu-3, the same base model under the same protocol, is healthy at \TuluGptqIF{} while its pair collapses its RTN network to \TuluPairIF{} (RTN \TuluRtnIF{}); and Qwen's window-calibrated sink matrix, from a protocol the model survives, collapses the RTN network to \QcwinLesionRtnIF{}.
So the lesion, with its layer-0 copy on Llama, collapses every RTN network we tried, where nothing is fitted around it, and a GPTQ network survives only when the matrices quantized after the lesion, solved on activations that already carry it, absorb it: on Llama the twelve other matrices of the first two layers, since rounding everything later to nearest changes nothing; on Qwen not yet localized.
Protocol and draw thus act on different links: short documents decide whether Llama's second sink exists, the draw whether its neighbors absorb the pair.
\Cref{sec:mechanism} applies to each lesion; what it does not give is the absorption.
